\section{Conclusions}
\label{sec:conclusion}

This paper presented a service-oriented 4D path-planning architecture for structured corridor operations within U-space. The architecture separates topology provision, A*-based corridor routing, waypoint synthesis, CFMS-based trajectory generation, NDMap-based deconfliction, and a USSP-facing authorisation boundary into distinct service responsibilities. The Planning Service owns route computation, trajectory generation, and request-local graph modification for lateral re-routing, while the Deconfliction Service owns NDMap conflict detection and conflict reporting. The USSP-facing boundary is downstream of these services: the planner prepares candidate 4D intent, but does not authorise U-space operations.

The Hamburg vertiport corridor-network evaluation produced two main results. First, in the planner-stage scaling experiment with deconfliction disabled, A* path computation and waypoint synthesis remained below one millisecond in the measured benchmark, while CFMS invocation and trajectory generation dominated planner-stage runtime and determined the observed scaling with route distance. Second, in the conflict-resolution-enabled experiment, the prototype retained 97 of 100 sequential requests as locally deconflicted candidate plans, including three requests resolved by lateral re-routing after initial conflicts. Only four requests had initial conflicts, so this result should be read as preliminary evidence for the feedback loop rather than as a statistical conflict-resolution-performance estimate. The three failed requests consisted of one retry-limit case and two routing/input-graph errors in the prototype re-planning loop.

These results provide evidence for the service-oriented prototype and identify CFMS trajectory generation as the main computational bottleneck, especially when conflict resolution requires repeated trajectory generation. They do not constitute an operational capacity assessment. Future work should evaluate larger and denser networks, concurrent request loads, distributed deployment effects, temporal and vertical deconfliction, and more realistic constraints including terrain, weather, obstacle avoidance, and dynamic traffic.
